\section{Measurable collision phases and symplectic area}
\label{sec:measurable-phase-rigidity}
The measurable coboundary in Theorem~\ref{thm:gaussian-kernel-coboundary}
need not have a value at a prescribed periodic orbit. We shall instead
use conditional pairs on a positive-measure hyperbolic product set.
The discontinuous section contribution will be removed algebraically
before this argument is applied. Throughout this section $R$ is fixed;
none of the local product constants is asserted uniform in $R$.

Write $\kappa_R$ for the one-collision lattice label, in the basis
$a=(1,0)$, $b=(1/2,\sqrt3/2)$. Thus
$f_R=(\kappa_{R,1},\kappa_{R,2},1,\tau_R)$. All identities between
measurable functions are understood almost everywhere with respect to
$\nu$. Let $\mathbb S^1=\{z\in\mathbb C:|z|=1\}$.

\subsection{A qualitative local product input}
\begin{lemma}[Regular product sets for the collision map]
\label{lem:regular-collision-product}
There is a compact positive-$\nu$-measure set $\mathcal P$ in a regular
collision coordinate rectangle, away from grazing, with the following
properties. It has product coordinates
$\Psi:K^u\times K^s\longrightarrow\mathcal P$. Write a product point as $\Psi(u,s)$. Fixing $u$ places points on one
local homogeneous stable curve; fixing $s$ places them on one local
homogeneous unstable curve. Thus $K^u$ labels stable leaves and $K^s$
labels unstable leaves. These curves are $C^1$, have transverse tangent cones,
and each stable curve meets each unstable curve exactly once.

The probability $m=\nu(\mathcal P)^{-1}\nu|_{\mathcal P}$ is equivalent,
in these coordinates, to a product $m^u\otimes m^s$ of nonatomic
transversal probabilities. A stable arc joining two points on the same
curve remains a connected regular arc under every forward collision
iterate, and its length tends to zero. The corresponding statement
holds for unstable arcs under backward iterates. In particular the
successive lattice labels agree along the respective arcs.

For every sufficiently small $\varepsilon>0$, the product coordinates
can be restricted to positive-measure subsets so that every four-corner
stable--unstable quadrilateral lies in a coordinate box of diameter
less than $\varepsilon$. Almost every such quadrilateral with distinct
coordinates is a simple, nondegenerate Jordan domain.
\end{lemma}
\begin{proof}
We use the qualitative billiard product construction in
\cite[Sections 8.2--8.3]{Young}. Section 8.2B gives transverse stable
and unstable cones. Homogeneous local curves and their regular iterates
are described in Section 8.2D and constructed in Section 8.3,
Sublemma 1. The product set chosen there has positive invariant area;
this is stated at the end of its construction, before the verification
of (P3)--(P5). The absolute-continuity formula (P5)(b), verified there,
has positive finite Jacobian. Apply the same local statement in the
reverse time direction. Since the invariant collision measure is
smooth and positive away from grazing, disintegration and the two
holonomy formulas identify its null sets on the product rectangle with
those of the product of the transversal measures. These measures are
nonatomic, being equivalent to restricted arc length on transversals.
Restriction to compact positive-measure subsets gives the stated form.
No assertion that $\mathcal P$ contains an open set is needed.

For completeness, contraction here means contraction of the actual
curves, not just their endpoints in a symbolic metric. The billiard
estimates of \cite[Section 8.2C]{Young} give exponential contraction of
the stable arcs in its adapted metric and
$\operatorname{length}(\gamma)\le C\sqrt{\operatorname{length}_{\rm ad}(\gamma)}$
for a stable or unstable curve. Thus, for fixed $R$ and the chosen product set, the ordinary lengths
of these iterated arcs are bounded by $C_{\mathcal P}\theta_{\mathcal P}^n$
for some $C_{\mathcal P}<\infty$ and $0<\theta_{\mathcal P}<1$.
In particular they tend to zero. No uniform choice in $R$ is made. Passing from $(\alpha,\varphi)$ to $(\alpha,p)$ uses only
$p=\sin\varphi$, whose derivative is bounded; no inverse change at
grazing is made. Homogeneous local curves do not cross the relevant
finite-iterate collision singularities. The label $\kappa_R$ is
locally constant on each such physical branch, which proves the label
assertion.

The hypotheses of that construction hold here: the disks are disjoint
with separation at least $3/50$, their boundaries are smooth with
strictly positive curvature, the free domain is connected, and the
horizon is finite. These are properties of the physical Euclidean
table. The same local construction works on its flat triangular torus;
alternatively the rectangular two-scatterer cover in
Lemma~\ref{lem:collision-spectral-input} gives the local product set in
a covering chart. There is no affine change of the reflection law.

Finally choose a product support point and restrict each transversal
to a sufficiently short positive-measure interval around it. Continuity
of the local product map and the transverse graph cones put the
connecting arcs into a small coordinate box. Unique intersections of
the graph curves make the four distinct corners bound a simple Jordan
domain. Nonatomicity excludes repeated transversal coordinates almost
everywhere. This also proves the last assertion for product sets with
gaps.
\end{proof}


\paragraph{Conditional pairs in product coordinates.}
To fix the measure conventions, write $\mu^u=m^u$, $\mu^s=m^s$,
$\mu=\mu^u\otimes\mu^s$ and
\[
 \dd m_{\mathcal P}=\rho(u,s)\dd\mu^u(u)\dd\mu^s(s),\qquad
 m_{\mathcal P}=\nu(\mathcal P)^{-1}\nu|_{\mathcal P}.
\]
Here $0<\rho<\infty$ almost everywhere and $\int\rho\dd\mu=1$.
Put $\rho_u(u)=\int\rho(u,s)\dd\mu^s(s)$ and
$\rho_s(s)=\int\rho(u,s)\dd\mu^u(u)$. The stable and unstable pair
probabilities are explicitly
\begin{align}\label{eq:explicit-conditional-pair-densities}
 \dd\mathfrak m_s(u,s_0,s_1)
 &=\frac{\rho(u,s_0)\rho(u,s_1)}{\rho_u(u)}
           \dd\mu^u(u)\dd\mu^s(s_0)\dd\mu^s(s_1),\\
 \dd\mathfrak m_u(u_0,u_1,s)
 &=\frac{\rho(u_0,s)\rho(u_1,s)}{\rho_s(s)}
           \dd\mu^u(u_0)\dd\mu^u(u_1)\dd\mu^s(s).\notag
\end{align}
The denominators are positive and finite almost everywhere. Integration
of either unused coordinate shows that both marginals are
$m_{\mathcal P}$. The displayed positive finite densities also give
exactly the triple-product measure equivalences used below, without
a boundedness assertion for $\rho$ or its reciprocal.

\subsection{The physical action one-form}
On the collision cylinder define
\begin{equation}\label{eq:collision-action-form}
 \vartheta_R=Rp\,\dd\alpha,\qquad
 \dd\vartheta_R=R\,\dd p\wedge\dd\alpha.
\end{equation}
The one-form is globally defined on the cylinder and bounded in both
the flat and the original collision coordinates.

\begin{lemma}[Exact action on a regular branch]
\label{lem:exact-collision-action}
On every regular collision branch,
\begin{equation}\label{eq:exact-collision-action}
 \dd\tau_R=T_R^*\vartheta_R-\vartheta_R.
\end{equation}
Consequently, if an arc $\gamma$ from $x$ to $y$ stays on regular
branches through time $n$, then
\begin{equation}\label{eq:action-telescope-arc}
 S_{n,R}\tau_R(y)-S_{n,R}\tau_R(x)
       =\int_{T_R^n\gamma}\vartheta_R-\int_\gamma\vartheta_R.
\end{equation}
For backward regular arcs the corresponding identity is
\begin{equation}\label{eq:backward-action-arc}
 S_{n,R}\tau_R(T_R^{-n}y)-S_{n,R}\tau_R(T_R^{-n}x)
       =\int_\gamma\vartheta_R-\int_{T_R^{-n}\gamma}\vartheta_R.
\end{equation}
\end{lemma}
\begin{proof}
Lift the flight to its actual initial disk and the next disk with fixed
label $d$. Its endpoints are $q=Rn_\alpha$ and
$q'=d+Rn_{\alpha'}$. If $v$ is its unit velocity, differentiation of
its Euclidean length gives
$\dd\tau_R=v\cdot\dd q'-v\cdot\dd q$.
The first tangential component is $p$, and specular reflection preserves
the tangential component at arrival, which is $p'$ in the outgoing
coordinates of $T_Rx$. Hence
$\dd\tau_R=Rp'\dd\alpha'-Rp\dd\alpha$. The lattice label is constant
on the branch and contributes no derivative. This proves
\eqref{eq:exact-collision-action}. Summing its pullbacks and integrating
over the arc proves \eqref{eq:action-telescope-arc}; applying that
identity to $T_R^{-n}\gamma$ proves \eqref{eq:backward-action-arc}.
\end{proof}

\begin{proposition}[A measurable physical phase has zero roof frequency]
\label{prop:measurable-roof-rigidity}
Suppose $q:M\to\mathbb S^1$ is measurable and, for constants
$u\in\R^2$ and $s,t\in\R$,
\begin{equation}\label{eq:physical-measurable-phase}
 q\circ T_R=\exp\{i(u\cdot\kappa_R+s+t\tau_R)\}\,q.
\end{equation}
Then $t=0$.
\end{proposition}
\begin{proof}
Discard an invariant null set so that every finite iterate of the phase
identity holds at the remaining points. Use the product set and the
probability $m$ of Lemma~\ref{lem:regular-collision-product}.
Disintegrate $m$ over its stable curves and, conditionally on the curve,
draw two independent points $x,y$. Denote this pair probability by
$\mathfrak m_s$. Both of its marginals are exactly $m$; in particular
they are bounded by $\nu(\mathcal P)^{-1}\nu$.

The following argument justifies endpoint agreement for a merely
measurable $q$. Let $q_\epsilon$ be a smooth approximation in $L^1(\nu)$,
with supremum norm at most one. Invariance of $\nu$ and the marginal
bound give, for every $n$,
\begin{align}\label{eq:conditional-pair-approximation}
 &\int |q(T_R^ny)-q(T_R^nx)|\dd\mathfrak m_s(x,y)\notag\\
 &\quad\le \frac{2}{\nu(\mathcal P)}\|q-q_\epsilon\|_{L^1(\nu)}
       +\int|q_\epsilon(T_R^ny)-q_\epsilon(T_R^nx)|\dd\mathfrak m_s(x,y).
\end{align}
For fixed $\epsilon$, the last term tends to zero by contraction and
bounded convergence. Then let $\epsilon$ tend to zero. Thus the left
side tends to zero. This uses bounded marginals, not invariance of the
pair probability or a claim about every stable curve.

Let $\gamma_s(x,y)$ be the stable arc oriented from $x$ to $y$.
Its forward labels agree at all times. Iteration of
\eqref{eq:physical-measurable-phase} and
\eqref{eq:action-telescope-arc} gives
\[
 \frac{q(T_R^ny)}{q(T_R^nx)}
 =\exp\!\left(it\left[\int_{T_R^n\gamma_s(x,y)}\vartheta_R
                         -\int_{\gamma_s(x,y)}\vartheta_R\right]\right)
                         \frac{q(y)}{q(x)}.
\]
The first arc integral tends to zero, because $\vartheta_R$ is bounded
and the arc length tends to zero. Since $|q|=1$, the left side minus one
has $L^1(\mathfrak m_s)$ norm tending to zero by
\eqref{eq:conditional-pair-approximation}. The right side has the
pointwise limit just displayed. Fatou's lemma therefore gives
\begin{equation}\label{eq:measurable-stable-holonomy}
 \frac{q(y)}{q(x)}
       =\exp\!\left(it\int_{\gamma_s(x,y)}\vartheta_R\right)
       \quad\text{for }\mathfrak m_s\text{-almost every }(x,y).
\end{equation}
Disintegrating instead over unstable curves and using backward iterates
in \eqref{eq:conditional-pair-approximation} gives, by
\eqref{eq:backward-action-arc},
\begin{equation}\label{eq:measurable-unstable-holonomy}
 \frac{q(y)}{q(x)}
       =\exp\!\left(it\int_{\gamma_u(x,y)}\vartheta_R\right)
       \quad\text{for almost every unstable pair}.
\end{equation}
The sign is the same in the two formulas: the backward action in
\eqref{eq:backward-action-arc} has limiting value
$\int_{\gamma_u(x,y)}\vartheta_R$.

We now use the measure-class part of the local product input. The stable
and unstable pair probabilities are equivalent, in product coordinates,
to the corresponding three-coordinate product measures. Fubini's theorem
therefore permits \eqref{eq:measurable-stable-holonomy} and
\eqref{eq:measurable-unstable-holonomy} to be multiplied around almost
every four-corner product quadrilateral $Q$. The oriented corner order is
\begin{gather*}
 x_{00}=\Psi(u_0,s_0)\ \longrightarrow\
 x_{01}=\Psi(u_0,s_1)\ \longrightarrow\ x_{11}=\Psi(u_1,s_1),\\
 x_{11}\ \longrightarrow\ x_{10}=\Psi(u_1,s_0)
            \ \longrightarrow\ x_{00}.
\end{gather*}
The edges are stable, unstable, stable with reversed $s$-orientation,
and unstable with reversed $u$-orientation. This specifies the boundary
orientation used in Stokes' formula; its sign need not be fixed when
only the absolute area is used. Their left sides telescope, so
\begin{equation}\label{eq:measurable-area-obstruction}
 \exp\!\left(it\oint_{\partial Q}\vartheta_R\right)=1.
\end{equation}
This assertion is on a positive-measure family of quadrilaterals; it does
not assign values to $q$ on a preselected closed orbit.

Suppose $t\ne0$. Restrict to a positive-measure product subrectangle
small enough that every such quadrilateral has coordinate area less
than $2\pi/(R|t|)$. Almost every quadrilateral is a nondegenerate Jordan
domain. Stokes' formula and \eqref{eq:collision-action-form} give
\[
 0<\left|t\oint_{\partial Q}\vartheta_R\right|
       =|t|R\operatorname{area}_{\alpha,p}(Q)<2\pi,
\]
contradicting \eqref{eq:measurable-area-obstruction}. Hence $t=0$.
\end{proof}

\subsection{Rotations and the constant collision phase}
Let $S_j(\alpha,p)=(\alpha+j\pi/3,p)$, $j\in\Z$. In lattice
coordinates rotation by $\pi/3$ is
\begin{equation}\label{eq:triangular-rotation-matrix}
 A=\begin{pmatrix}0&-1\\1&1\end{pmatrix},\qquad
 A^3=-I_2,\qquad I_2+A^2+A^4=0.
\end{equation}
The actual physical symmetries give
\begin{equation}\label{eq:collision-rotation-identities}
 T_RS_j=S_jT_R,\qquad
 \kappa_R\circ S_j=A^j\kappa_R,\qquad
 \tau_R\circ S_j=\tau_R,\qquad (S_j)_*\nu=\nu.
\end{equation}
They follow by rotating both endpoints, the velocity and the selected
next disk. The first admissible physical root remains the first root.
No invariance of $Y_R^*$ under these rotations is asserted.

\begin{theorem}[Rigidity of the roof and constant parts of a phase]
\label{thm:physical-phase-rigidity}
If a measurable circle-valued $q$ satisfies
\eqref{eq:physical-measurable-phase}, then
\begin{equation}\label{eq:physical-phase-conclusion}
 t=0,\qquad s\in2\pi\Z.
\end{equation}
\end{theorem}
\begin{proof}
The proposition gives $t=0$. Define
\[
 Q_2=q\,(q\circ S_3),\qquad
 Q_3=q\,(q\circ S_2)(q\circ S_4).
\]
By \eqref{eq:triangular-rotation-matrix} and
\eqref{eq:collision-rotation-identities}, the lattice phases cancel
exactly and
\[
 Q_2\circ T_R=e^{2is}Q_2,\qquad Q_3\circ T_R=e^{3is}Q_3.
\]
Both functions have modulus one. The mixing collision map has no
nontrivial circle eigenvalue, by the $L^2$ argument in the proof of
Corollary~\ref{cor:positive-count-variance}. Therefore
$e^{2is}=e^{3is}=1$, and their quotient gives $e^{is}=1$.
\end{proof}

The theorem is qualitative. It excludes exact measurable physical
phases with nonzero roof frequency or a nonintegral constant phase.
It is not a norm estimate for approximate spectral vectors, and supplies
no decay rate for a high-frequency collision or induced resolvent.
